\section*{Chapter \ref{ch:g7:finding-and-fixing-bugs} --- Finding and Fixing Bugs}

\begin{solution}{exo:g7:finding-and-fixing-bugs:1}
(a) \omterm{def:g7:finding-and-fixing-bugs:kinds}{syntax error}; (b) \omterm{def:g7:finding-and-fixing-bugs:kinds}{runtime error}; (c) \omterm{def:g7:finding-and-fixing-bugs:kinds}{logic error}; (d) \omterm{def:g7:finding-and-fixing-bugs:kinds}{runtime error}: the
\omterm{def:g4:programs-and-operating-system:program}{program} starts, and stops with a \code{NameError} when it reaches that line;
(e) \omterm{def:g7:finding-and-fixing-bugs:kinds}{logic error}; (f) \omterm{def:g7:finding-and-fixing-bugs:kinds}{syntax error}.
\end{solution}

\begin{solution}{exo:g7:finding-and-fixing-bugs:2}
From the bottom: an \code{IndexError}, a \omterm{def:g7:lists:list}{list} \omterm{def:g7:lists:list}{index} out of range, on line 2,
inside the \omterm{def:g7:functions:function}{function} \code{last}, called from line 5. A \omterm{def:g7:lists:list}{list} of 3 \omterm{def:g7:lists:list}{items} has
indices 0 to 2, and \code{values[len(values)]} asks for \omterm{def:g7:lists:list}{index} 3; the last
\omterm{def:g7:lists:list}{item} is \code{values[len(values) - 1]}, or \code{values[-1]}.
\end{solution}

\begin{solution}{exo:g7:finding-and-fixing-bugs:3}
The six lines that start with \code{trace:}. Only the last line, 13.25,
remains.
\end{solution}

\begin{solution}{exo:g7:finding-and-fixing-bugs:4}
Each opening bracket needs its closing bracket; \code{print} opens one and
\code{len} opens another, so two are needed at the end:
\code{print("count:", len(marks))}.
\end{solution}

\begin{solution}{exo:g7:finding-and-fixing-bugs:5}
\code{range(6)} gives the indices 0 to 5: six marks, the sixth being 14
($53 + 14 = 67$). Write \code{range(5)}, and the \omterm{def:g4:programs-and-operating-system:program}{program} prints 53.
\end{solution}

\begin{solution}{exo:g7:finding-and-fixing-bugs:6}
The trace has one line only: the \omterm{def:g4:algorithms-repeat-and-decide:loop}{loop} stopped after its first \omterm{def:g6:variables-and-input:value}{value}. The
\code{return} is indented inside the \omterm{def:g4:algorithms-repeat-and-decide:loop}{loop}, so the \omterm{def:g7:functions:function}{function} returns at the
end of the first turn. Moved back to the level of \code{for}, after the \omterm{def:g4:algorithms-repeat-and-decide:loop}{loop},
it gives 2.
\end{solution}

\begin{solution}{exo:g7:finding-and-fixing-bugs:7}
The line \code{print("  trace: end", end)}, just after \code{end} is computed,
prints 3, 5, $-1$ and $-1$; for \code{"Stop"} it prints $-1$, which shows that
\code{find} found no space.
\end{solution}

\begin{solution}{exo:g7:finding-and-fixing-bugs:8}
A \omterm{def:g7:finding-and-fixing-bugs:kinds}{syntax error}: Python reports \code{invalid syntax} and runs nothing. The
\omterm{def:g6:if-else:comparison}{comparison} is written \code{<=}, the \code{<} first, as it is said: ``less
than or equal to''.
\end{solution}

\begin{solution}{exo:g7:finding-and-fixing-bugs:9}
\code{values[i] < values[i + 1]} is false for two equal neighbours, so a \omterm{def:g7:lists:list}{list}
with equal neighbours is called not in order. The smallest such \omterm{def:g7:lists:list}{list} is
\code{[1, 1]}. The \omterm{def:g6:if-else:comparison}{comparison} must be \code{<=}:

\omcode{code/grade-7/05-finding-and-fixing-bugs/exercises/is_sorted_fixed.py}{1}{10}{Equal
neighbours allowed.}\label{lst:g7:finding-and-fixing-bugs:sol-is-sorted}

It prints \code{True False} and \code{True True}.
\end{solution}

\begin{solution}{exo:g7:finding-and-fixing-bugs:10}
A test before the division, which is never reached for an empty \omterm{def:g7:lists:list}{list}:

\omcode{code/grade-7/05-finding-and-fixing-bugs/exercises/average_guard.py}{1}{10}{An
average with a guard.}\label{lst:g7:finding-and-fixing-bugs:sol-guard}

It prints \code{10.6 0}.
\end{solution}

\begin{solution}{exo:g7:finding-and-fixing-bugs:11}
The two \omterm{def:g7:finding-and-fixing-bugs:bug}{bugs}: \code{range(1, n)} stops at 9, so it must be
\code{range(1, n + 1)}; and \texttt{i \% 2 == 1} picks the odd numbers, so it
must be \texttt{i \% 2 == 0}. Fixing the test first gives $2 + 4 + 6 + 8 = 20$,
then 30; fixing the range first gives $1 + 3 + 5 + 7 + 9 = 25$ again, then 30.
Both orders end with the same \omterm{def:g4:programs-and-operating-system:program}{program}, which prints 30:

\omcode{code/grade-7/05-finding-and-fixing-bugs/exercises/evens_fixed.py}{1}{6}{Both
\omterm{def:g7:finding-and-fixing-bugs:bug}{bugs} fixed.}\label{lst:g7:finding-and-fixing-bugs:sol-evens}

Running after each fix tells you what each change did: in the second order,
the first fix leaves the answer at 25, and only running shows that it was
still needed. With both changes at once, a wrong result would not say which
change was wrong.
\end{solution}

\begin{solution}{exo:g7:finding-and-fixing-bugs:12}
The \omterm{def:g4:algorithms-repeat-and-decide:loop}{loop} starts at \omterm{def:g7:lists:list}{index} 1, so the first mark is never looked at, and
\code{best} starts at 0. The \omterm{def:g7:functions:function}{function} is wrong when the largest mark is the
first one: \code{[9, 3, 4]} gives 4. The smallest such \omterm{def:g7:lists:list}{list} is a \omterm{def:g7:lists:list}{list} of one
mark, \code{[5]}, which gives 0. Starting with \code{best = values[0]} fixes
both: the \omterm{def:g4:programs-and-operating-system:program}{program} then prints \code{9 9 5} for \code{[3, 9, 4]}, \code{[9, 3,
4]} and \code{[5]}.
\end{solution}

\begin{solution}{pb:g7:finding-and-fixing-bugs:1}
\textbf{1.} Line 17, \code{def comment(score)}, has no colon at the end.

\textbf{2.} The first report points to line 17, the missing colon; the
second, read from the bottom, to line 19, \code{return WORDS[tens]}.

\textbf{3.} \code{comment} failed, called by \code{grade} for the score 99.
\code{tens} was 9, and \code{WORDS} has only nine \omterm{def:g7:lists:list}{items}, with indices 0 to 8:
the word \code{"great"} is missing. Adding it fixes the \omterm{def:g7:finding-and-fixing-bugs:kinds}{runtime error}.

\textbf{4.} Answers typed by hand can be mistyped, and there would be a
hundred of them; a second \omterm{def:g4:programs-and-operating-system:program}{program} written another way is unlikely to make
the same mistake as the first.

\textbf{5.}

\omoutput{code/grade-7/05-finding-and-fixing-bugs/outputs/check_grader.txt}{66,
then 1, then 0 wrong answers.}\label{out:g7:finding-and-fixing-bugs:check}

\textbf{6.} Between \code{v2} and \code{v3}, \texttt{score \% 10} became
\code{score // 10}: the \omterm{def:g6:first-python-program:comment}{comment} used the units digit instead of the tens
digit. Between \code{v3} and \code{v4}, \code{score > 90} became
\code{score >= 90}.

\textbf{7.} Yes: 45 should be ``poor'' and 75 ``good''. Their units digit is
5, and \code{WORDS[5]} is ``fair''.

\textbf{8.} The score 5: for 0 to 4 the units digit and the tens digit both
pick ``poor'', and 5 is the first score whose units digit picks another word.

\textbf{9.} Only 90, which gets a B. Part I's \omterm{def:g1:what-is-a-computer:input}{output} already showed it: the
line \code{90 B (poor)}.

\textbf{10.} 66 still: the score 90 also has a wrong \omterm{def:g6:first-python-program:comment}{comment}, so fixing its
letter alone does not make its answer right.

\textbf{11.} \textbf{Result.} Among the scores 0 to 99, the grader gives 66
wrong answers once it runs, 1 after the \omterm{def:g6:first-python-program:comment}{comment} is fixed and 0 after the
letter is fixed.
\end{solution}
